% !PRINTED 237
\weberpage{237}
%<<CONT
Sind $a$, $b$, $c$, $\partial$ vier der Bedingung:
\[
a\partial - bc = n
\]
genügende ganze Zahlen ohne gemeinsamen Teiler, so erhält man
die einer Reihe entsprechenden Zahlenpaare $\mu$, $\mu'$, wenn man in
\weq{19}{\begin{gathered}
\mu \equiv a\alpha' + c\beta' \\
\mu' \equiv b\alpha' + \partial\beta' \ (\mathrm{mod}\ n)
\end{gathered}}
$\alpha'$, $\beta'$ alle und nur solche Werte durchlaufen läßt, bei denen
$\mu$, $\mu'$ ohne gemeinsamen Teiler mit $n$ sind.

Daß zwei den Kongruenzen (19) entsprechende Wertpaare $\mu$, $\mu'$
wirklich derselben Reihe angehören, ergibt sich unmittelbar
aus § 65 (3), und ebenso ist selbstverständlich, daß \gesperrt{alle Zahlenpaare}
einer Reihe in dieser Form enthalten sind, da man $\alpha'$, $\beta'$
durch $h\alpha'$, $h\beta'$ ersetzen kann, wenn $h$ relativ prim zu $n$ ist. Daß
man \gesperrt{sämtliche Reihen} auf diesem Wege bekommt, zeigen
die Formeln (10).

\begin{center}
\textbf{§ 69.} \textbf{\gesperrt{Die Invariantengleichung.}}
\end{center}

Unter den Transformationsgleichungen verdient eine ein besonderes
Interesse, nämlich die, deren Wurzeln die $\psi(n)$ Größen
\weq{1}{j\left(\frac{c+\partial\omega}{a}\right)}
oder nach der Bestimmung (19), § 68 die damit identischen
Größen
\weq{2}{j\left(\frac{c+\partial\omega}{a+b\omega}\right)}
sind, wenn $j(\omega)$ die in § 46 definierte Invariante ist.

Diese Gleichung heißt die \gesperrt{Invariantengleichung.}

Die Funktion $j(\omega)$ läßt sich nach § 54 rational durch $f(\omega)^{24}$
darstellen, nämlich
\weq{3}{j(\omega) = \frac{[f(\omega)^{24} - 16]^3}{f(\omega)^{24}},}
so daß also nach den Resultaten des vorigen Paragraphen die
Größen (1) die Wurzeln einer Gleichung sind, deren Koeffizienten
rationale Funktionen von $\varkappa^2$ sind.

Setzt man aber für $f(\omega)$ in (3) eine der früher gefundenen
Entwickelungen, z. B. § 24, (11):
\[
% [sic] The source prints a stray comma after the lower limit 1.
f(\omega) = q^{-\frac{1}{24}} \prod_{1,}^{\nu} (1+q^{2\nu-1}),
\]
%CONT>>
